\chapter{高频算法模板库}

\section{模板库使用方法}

本章把前面章节中反复出现的算法模板统一编号。考场上建议先通过决策树确定题型，再来本章找对应模板。

\begin{itemize}
  \item \Code{T001--T099}：基础实现和 STL。
  \item \Code{T100--T199}：前缀和、差分、排序、哈希、二分。
  \item \Code{T200--T299}：搜索和图论。
  \item \Code{T300--T399}：动态规划。
  \item \Code{T400--T499}：数据结构。
  \item \Code{T500--T599}：字符串、数学、日期。
\end{itemize}

\section{模板速查表}

\begin{longtable}{p{2cm}p{4cm}p{5cm}p{3cm}}
\toprule
编号 & 名称 & 识别信号 & 复杂度 \\
\midrule
T101 & 一维前缀和 & 多次静态区间和查询 & 预处理 $O(n)$，查询 $O(1)$ \\
T102 & 二维前缀和 & 多次子矩阵和查询 & 预处理 $O(nm)$，查询 $O(1)$ \\
T103 & 一维差分 & 多次区间加，最后查询 & $O(n+q)$ \\
T104 & 二维差分 & 多次矩形加，最后输出 & $O(nm+q)$ \\
T105 & 排序自定义比较器 & 多关键字排序、排名 & $O(n\log n)$ \\
T106 & 哈希统计 & 出现次数、查重、频率 & 平均 $O(n)$ \\
T107 & 二分答案 & 最大化最小值、最小化最大值 & $O(\log V \cdot check)$ \\
T108 & 双指针 & 排序后找对、滑动窗口 & $O(n)$ 或 $O(n\log n)$ \\
T109 & 坐标离散化 & 坐标大但点数少 & $O(n\log n)$ \\
T110 & 事件排序 & 按时间发生、扫描线 & $O(n\log n)$ \\
T201 & 网格 BFS & 最短步数、迷宫 & $O(nm)$ \\
T202 & DFS 连通块 & 区域大小、连通块数量 & $O(n+m)$ 或 $O(nm)$ \\
T203 & 并查集 & 合并集合、判断同组 & 近似 $O(1)$ \\
T204 & Dijkstra & 正权最短路 & $O((n+m)\log n)$ \\
T205 & 拓扑排序 & 依赖关系、先后顺序 & $O(n+m)$ \\
T206 & Floyd & 小图任意两点最短路 & $O(n^3)$ \\
T207 & Kruskal & 最小生成树 & $O(m\log m)$ \\
T208 & LCA 倍增 & 树上最近公共祖先 & $O(n\log n)$ 预处理，$O(\log n)$ 查询 \\
T301 & 01 背包 & 每个物品选或不选 & $O(nV)$ \\
T302 & 完全背包 & 每个物品可选多次 & $O(nV)$ \\
T401 & 树状数组 & 单点修改、区间和 & $O(\log n)$ \\
T402 & 线段树 & 区间修改、区间查询 & $O(\log n)$ \\
T403 & 单调栈 & 最近更大/更小元素 & $O(n)$ \\
T404 & 单调队列 & 滑动窗口最大/最小 & $O(n)$ \\
T501 & split 字符串分割 & 路径、配置、命令解析 & $O(|s|)$ \\
T502 & KMP & 长文本匹配模式串 & $O(n+m)$ \\
T503 & Trie & 字符串集合、前缀查询 & $O(\sum |s|)$ \\
T510 & 整除取整和标准取模 & 分组数、负数取模、数学 floor/ceil & $O(1)$ \\
T511 & 快速幂 & $a^b \bmod m$ & $O(\log b)$ \\
T512 & 素数和约数 & 判断素数、枚举约数 & $O(\sqrt n)$ \\
T513 & 进制转换 & 任意进制与十进制 & $O(\text{位数})$ \\
T514 & 日期处理 & 闰年、每月天数、日期推进 & 按实现而定 \\
T601 & 大模拟骨架 & 多对象、多命令、长题面 & 按题目而定 \\
T602 & 命令分发 & 每行一个操作名 & 按命令复杂度 \\
T603 & 状态机 & 对象多状态切换 & 按操作复杂度 \\
T604 & 懒删除优先队列 & 动态删除/更新后仍取最优 & $O(\log n)$ 均摊 \\
T605 & LRU 缓存 & 最近最少使用淘汰 & $O(1)$ 均摊 \\
T606 & 路径规范化 & \Code{/a/../b}、目录路径 & $O(|s|)$ \\
T607 & URL 参数解析 & \Code{path?a=1\&b=2} & $O(|s|)$ \\
T608 & 括号匹配 & 括号合法性、嵌套结构 & $O(n)$ \\
T609 & 简单表达式求值 & 只含加减的表达式 & $O(n)$ \\
\bottomrule
\end{longtable}

\section{T101 一维前缀和}

\subsection{适用场景}

\begin{itemize}
  \item 数组不修改。
  \item 多次查询 $[l,r]$ 的和。
  \item 需要快速得到一段累计值。
\end{itemize}

\subsection{模板}

\begin{lstlisting}
int n, q;
cin >> n >> q;
vector<long long> pre(n + 1, 0);

for (int i = 1; i <= n; i++) {
    long long x;
    cin >> x;
    pre[i] = pre[i - 1] + x;
}

while (q--) {
    int l, r;
    cin >> l >> r;
    cout << pre[r] - pre[l - 1] << '\n';
}
\end{lstlisting}

\subsection{易错点}

\begin{itemize}
  \item 模板默认下标从 1 开始。
  \item \Code{pre[0]} 是 0。
  \item 求和用 \Code{long long}。
  \item 如果有修改，普通前缀和不适用。
\end{itemize}

\section{T102 二维前缀和}

\subsection{适用场景}

多次查询子矩阵和。

\begin{lstlisting}
int n, m, q;
cin >> n >> m >> q;
vector<vector<long long>> pre(n + 1, vector<long long>(m + 1, 0));

for (int i = 1; i <= n; i++) {
    for (int j = 1; j <= m; j++) {
        long long x;
        cin >> x;
        pre[i][j] = pre[i - 1][j] + pre[i][j - 1]
                  - pre[i - 1][j - 1] + x;
    }
}

while (q--) {
    int x1, y1, x2, y2;
    cin >> x1 >> y1 >> x2 >> y2;
    long long ans = pre[x2][y2] - pre[x1 - 1][y2]
                  - pre[x2][y1 - 1] + pre[x1 - 1][y1 - 1];
    cout << ans << '\n';
}
\end{lstlisting}

\subsection{易错点}

\begin{itemize}
  \item 四项公式不要写反。
  \item 默认坐标从 1 开始。
  \item 子矩阵端点是否保证 $x1 \le x2, y1 \le y2$。
\end{itemize}

\section{T103 一维差分}

\subsection{适用场景}

多次区间加，最后统一输出每个位置的值。

\begin{lstlisting}
int n, q;
cin >> n >> q;
vector<long long> diff(n + 2, 0);

while (q--) {
    int l, r;
    long long v;
    cin >> l >> r >> v;
    diff[l] += v;
    diff[r + 1] -= v;
}

long long cur = 0;
for (int i = 1; i <= n; i++) {
    cur += diff[i];
    cout << cur << '\n';
}
\end{lstlisting}

\subsection{易错点}

\begin{itemize}
  \item 数组开到 \Code{n+2}。
  \item 如果修改后立刻查询区间和，普通差分不够。
  \item 修改值用 \Code{long long}。
\end{itemize}

\section{T104 二维差分}

\subsection{适用场景}

多次对子矩形加值，最后输出整个矩阵。

\begin{lstlisting}
int n, m, q;
cin >> n >> m >> q;
vector<vector<long long>> diff(n + 2, vector<long long>(m + 2, 0));

while (q--) {
    int x1, y1, x2, y2;
    long long v;
    cin >> x1 >> y1 >> x2 >> y2 >> v;
    diff[x1][y1] += v;
    diff[x2 + 1][y1] -= v;
    diff[x1][y2 + 1] -= v;
    diff[x2 + 1][y2 + 1] += v;
}

for (int i = 1; i <= n; i++) {
    for (int j = 1; j <= m; j++) {
        diff[i][j] += diff[i - 1][j] + diff[i][j - 1]
                    - diff[i - 1][j - 1];
    }
}
\end{lstlisting}

\section{T105 排序自定义比较器}

\subsection{适用场景}

多关键字排序、排名、优先级排序。

\begin{lstlisting}
struct Node {
    int id;
    int score;
    int time;
};

sort(a.begin(), a.end(), [](const Node& x, const Node& y) {
    if (x.score != y.score) return x.score > y.score;
    if (x.time != y.time) return x.time < y.time;
    return x.id < y.id;
});
\end{lstlisting}

\subsection{易错点}

\begin{itemize}
  \item 相等时不能返回 \Code{true}。
  \item 每个关键字升序还是降序要写清楚。
  \item 需要保持原顺序时用 \Code{stable\_sort}。
\end{itemize}

\section{T106 哈希统计}

\subsection{值域小}

\begin{lstlisting}
vector<int> cnt(100001, 0);
for (int x : a) {
    cnt[x]++;
}
\end{lstlisting}

\subsection{值域大或字符串}

\begin{lstlisting}
unordered_map<string, int> cnt;
for (string s : words) {
    cnt[s]++;
}
\end{lstlisting}

\subsection{易错点}

\begin{itemize}
  \item 需要有序输出时用 \Code{map}。
  \item key 很大不能开计数数组。
  \item \Code{cnt[x]} 会自动创建 key。
\end{itemize}

\section{T107 二分答案}

\subsection{适用场景}

\begin{itemize}
  \item 求最小可行值。
  \item 求最大可行值。
  \item 给定答案 $mid$ 能判断是否可行。
  \item 可行性具有单调性。
\end{itemize}

\subsection{最小可行值}

\begin{lstlisting}
bool check(long long mid) {
    // return true if mid is feasible
}

long long l = 0, r = 1000000000000000000LL;
while (l < r) {
    long long mid = l + (r - l) / 2;
    if (check(mid)) r = mid;
    else l = mid + 1;
}
cout << l << '\n';
\end{lstlisting}

\subsection{最大可行值}

\begin{lstlisting}
long long l = 0, r = 1000000000000000000LL;
while (l < r) {
    long long mid = l + (r - l + 1) / 2;
    if (check(mid)) l = mid;
    else r = mid - 1;
}
cout << l << '\n';
\end{lstlisting}

\Warn{二分答案必须有单调性。}

\section{T108 双指针}

\subsection{排序后找两数和}

\begin{lstlisting}
sort(a.begin(), a.end());
int l = 0, r = n - 1;
bool ok = false;

while (l < r) {
    long long sum = 1LL * a[l] + a[r];
    if (sum == target) {
        ok = true;
        break;
    } else if (sum < target) {
        l++;
    } else {
        r--;
    }
}
\end{lstlisting}

\subsection{非负数组滑动窗口}

\begin{lstlisting}
int l = 0;
long long sum = 0;
int best = 0;

for (int r = 0; r < n; r++) {
    sum += a[r];
    while (sum > limit && l <= r) {
        sum -= a[l];
        l++;
    }
    best = max(best, r - l + 1);
}
\end{lstlisting}

\Warn{滑动窗口通常要求元素非负或条件具有单调性。}

\section{T109 坐标离散化}

\subsection{适用场景}

坐标值很大，但出现的坐标数量不多。

\begin{lstlisting}
vector<int> xs;
for (int x : original) {
    xs.push_back(x);
}

sort(xs.begin(), xs.end());
xs.erase(unique(xs.begin(), xs.end()), xs.end());

auto getId = [&](int x) {
    return (int)(lower_bound(xs.begin(), xs.end(), x) - xs.begin()) + 1;
};
\end{lstlisting}

\Warn{如果题目涉及区间长度，不能只看离散编号，还要保留真实坐标差。}

\section{T110 事件排序}

\subsection{适用场景}

按时间顺序处理事件，或扫描区间开始/结束。

\begin{lstlisting}
struct Event {
    int time;
    int type;
    int id;
};

sort(events.begin(), events.end(), [](const Event& a, const Event& b) {
    if (a.time != b.time) return a.time < b.time;
    return a.type < b.type;
});

for (const Event& e : events) {
    if (e.type == 0) {
        // start
    } else {
        // end
    }
}
\end{lstlisting}

\subsection{区间同时进行数量}

\begin{lstlisting}
vector<pair<int, int>> events;
for (auto [l, r] : intervals) {
    events.push_back({l, 1});
    events.push_back({r, -1});
}

sort(events.begin(), events.end(), [](auto a, auto b) {
    if (a.first != b.first) return a.first < b.first;
    return a.second < b.second; // end before start for [l,r)
});

int cur = 0, best = 0;
for (auto [t, delta] : events) {
    cur += delta;
    best = max(best, cur);
}
\end{lstlisting}

\Warn{同一时间开始和结束谁先处理，要按题意决定。}

\section{T201 网格 BFS}

\subsection{适用场景}

迷宫、最短步数、每一步代价相同。

\begin{lstlisting}
int n, m;
cin >> n >> m;
vector<string> g(n);
for (int i = 0; i < n; i++) cin >> g[i];

int sx, sy, tx, ty;
cin >> sx >> sy >> tx >> ty;

vector<vector<int>> dist(n, vector<int>(m, -1));
queue<pair<int, int>> q;

dist[sx][sy] = 0;
q.push({sx, sy});

int dx[4] = {-1, 0, 1, 0};
int dy[4] = {0, 1, 0, -1};

while (!q.empty()) {
    auto [x, y] = q.front();
    q.pop();

    for (int d = 0; d < 4; d++) {
        int nx = x + dx[d];
        int ny = y + dy[d];
        if (nx < 0 || nx >= n || ny < 0 || ny >= m) continue;
        if (g[nx][ny] == '#') continue;
        if (dist[nx][ny] != -1) continue;
        dist[nx][ny] = dist[x][y] + 1;
        q.push({nx, ny});
    }
}

cout << dist[tx][ty] << '\n';
\end{lstlisting}

\subsection{易错点}

\begin{itemize}
  \item 坐标是 0 下标还是 1 下标。
  \item 入队时标记访问。
  \item 边权不全相等时不要用普通 BFS。
\end{itemize}

\section{T202 DFS 连通块}

\subsection{网格 DFS}

\begin{lstlisting}
int n, m;
vector<string> g;
vector<vector<int>> vis;
int dx[4] = {-1, 0, 1, 0};
int dy[4] = {0, 1, 0, -1};

int dfs(int x, int y) {
    vis[x][y] = 1;
    int area = 1;
    for (int d = 0; d < 4; d++) {
        int nx = x + dx[d];
        int ny = y + dy[d];
        if (nx < 0 || nx >= n || ny < 0 || ny >= m) continue;
        if (vis[nx][ny]) continue;
        if (g[nx][ny] == '#') continue;
        area += dfs(nx, ny);
    }
    return area;
}
\end{lstlisting}

\subsection{易错点}

\begin{itemize}
  \item 大网格递归可能爆栈，可改 BFS。
  \item 障碍字符要按题目改。
\end{itemize}

\section{T203 并查集}

\begin{lstlisting}
struct DSU {
    vector<int> parent, sz;

    DSU(int n) {
        parent.resize(n + 1);
        sz.assign(n + 1, 1);
        for (int i = 1; i <= n; i++) parent[i] = i;
    }

    int find(int x) {
        if (parent[x] == x) return x;
        return parent[x] = find(parent[x]);
    }

    bool unite(int a, int b) {
        int ra = find(a);
        int rb = find(b);
        if (ra == rb) return false;
        if (sz[ra] < sz[rb]) swap(ra, rb);
        parent[rb] = ra;
        sz[ra] += sz[rb];
        return true;
    }

    bool same(int a, int b) {
        return find(a) == find(b);
    }
};
\end{lstlisting}

\subsection{适用场景}

合并集合、判断是否同组、连通块。

\section{T204 Dijkstra}

\begin{lstlisting}
struct Edge {
    int to;
    long long w;
};

const long long INF = (long long)4e18;

int n, m, s;
cin >> n >> m >> s;
vector<vector<Edge>> g(n + 1);

for (int i = 0; i < m; i++) {
    int u, v;
    long long w;
    cin >> u >> v >> w;
    g[u].push_back({v, w});
    g[v].push_back({u, w}); // remove if directed
}

vector<long long> dist(n + 1, INF);
priority_queue<pair<long long, int>,
               vector<pair<long long, int>>,
               greater<pair<long long, int>>> pq;

dist[s] = 0;
pq.push({0, s});

while (!pq.empty()) {
    auto [d, u] = pq.top();
    pq.pop();
    if (d != dist[u]) continue;

    for (auto e : g[u]) {
        if (dist[e.to] > dist[u] + e.w) {
            dist[e.to] = dist[u] + e.w;
            pq.push({dist[e.to], e.to});
        }
    }
}
\end{lstlisting}

\subsection{易错点}

\begin{itemize}
  \item 不能处理负边权。
  \item 无向图双向加边，有向图单向加边。
  \item 距离用 \Code{long long}。
\end{itemize}

\section{T205 拓扑排序}

\subsection{适用场景}

任务依赖、课程先修、判断有向图是否有环。

\begin{lstlisting}
int n, m;
cin >> n >> m;
vector<vector<int>> g(n + 1);
vector<int> indeg(n + 1, 0);

for (int i = 0; i < m; i++) {
    int u, v;
    cin >> u >> v; // u before v
    g[u].push_back(v);
    indeg[v]++;
}

queue<int> q;
for (int i = 1; i <= n; i++) {
    if (indeg[i] == 0) q.push(i);
}

vector<int> order;
while (!q.empty()) {
    int u = q.front();
    q.pop();
    order.push_back(u);
    for (int v : g[u]) {
        indeg[v]--;
        if (indeg[v] == 0) q.push(v);
    }
}

if ((int)order.size() == n) {
    // has topological order
} else {
    // has cycle
}
\end{lstlisting}

\section{T206 Floyd}

\subsection{适用场景}

点数小，任意两点最短路。

\begin{lstlisting}
const long long INF = (long long)4e18;

int n, m;
cin >> n >> m;
vector<vector<long long>> dist(n + 1, vector<long long>(n + 1, INF));

for (int i = 1; i <= n; i++) dist[i][i] = 0;

for (int i = 0; i < m; i++) {
    int u, v;
    long long w;
    cin >> u >> v >> w;
    dist[u][v] = min(dist[u][v], w);
    dist[v][u] = min(dist[v][u], w); // remove if directed
}

for (int k = 1; k <= n; k++) {
    for (int i = 1; i <= n; i++) {
        for (int j = 1; j <= n; j++) {
            if (dist[i][k] == INF || dist[k][j] == INF) continue;
            dist[i][j] = min(dist[i][j], dist[i][k] + dist[k][j]);
        }
    }
}
\end{lstlisting}

\Warn{三层循环顺序必须是 \Code{k,i,j}。}

\section{T207 Kruskal}

\subsection{适用场景}

无向图，连接所有点的最小总代价。

\begin{lstlisting}
struct Edge {
    int u, v;
    long long w;
};

vector<Edge> edges;
sort(edges.begin(), edges.end(), [](const Edge& a, const Edge& b) {
    return a.w < b.w;
});

DSU dsu(n);
long long ans = 0;
int cnt = 0;

for (auto e : edges) {
    if (dsu.unite(e.u, e.v)) {
        ans += e.w;
        cnt++;
    }
}

if (cnt == n - 1) cout << ans << '\n';
else cout << "Impossible\n";
\end{lstlisting}

\section{T208 LCA 倍增}

\subsection{适用场景}

树上最近公共祖先、树上路径查询。

\begin{lstlisting}
const int LOG = 20;
vector<vector<int>> up;
vector<int> depth;
vector<vector<int>> g;

void dfsLca(int u, int p) {
    up[0][u] = p;
    for (int k = 1; k < LOG; k++) {
        up[k][u] = up[k - 1][up[k - 1][u]];
    }
    for (int v : g[u]) {
        if (v == p) continue;
        depth[v] = depth[u] + 1;
        dfsLca(v, u);
    }
}

int lca(int a, int b) {
    if (depth[a] < depth[b]) swap(a, b);
    int diff = depth[a] - depth[b];
    for (int k = 0; k < LOG; k++) {
        if (diff & (1 << k)) a = up[k][a];
    }
    if (a == b) return a;
    for (int k = LOG - 1; k >= 0; k--) {
        if (up[k][a] != up[k][b]) {
            a = up[k][a];
            b = up[k][b];
        }
    }
    return up[0][a];
}
\end{lstlisting}

初始化：

\begin{lstlisting}
up.assign(LOG, vector<int>(n + 1));
depth.assign(n + 1, 0);
dfsLca(1, 1);
\end{lstlisting}

\section{T301 01 背包}

\begin{lstlisting}
int n, V;
cin >> n >> V;
vector<int> w(n + 1), val(n + 1);
for (int i = 1; i <= n; i++) {
    cin >> w[i] >> val[i];
}

vector<int> dp(V + 1, 0);
for (int i = 1; i <= n; i++) {
    for (int j = V; j >= w[i]; j--) {
        dp[j] = max(dp[j], dp[j - w[i]] + val[i]);
    }
}
cout << dp[V] << '\n';
\end{lstlisting}

\Warn{01 背包容量倒序。}

\section{T302 完全背包}

\begin{lstlisting}
int n, V;
cin >> n >> V;
vector<int> w(n + 1), val(n + 1);
for (int i = 1; i <= n; i++) {
    cin >> w[i] >> val[i];
}

vector<int> dp(V + 1, 0);
for (int i = 1; i <= n; i++) {
    for (int j = w[i]; j <= V; j++) {
        dp[j] = max(dp[j], dp[j - w[i]] + val[i]);
    }
}
cout << dp[V] << '\n';
\end{lstlisting}

\Warn{完全背包容量正序。}

\section{T401 树状数组}

\begin{lstlisting}
struct BIT {
    int n;
    vector<long long> tree;

    BIT(int n) : n(n), tree(n + 1, 0) {}

    void add(int idx, long long val) {
        for (; idx <= n; idx += idx & -idx) {
            tree[idx] += val;
        }
    }

    long long sumPrefix(int idx) {
        long long res = 0;
        for (; idx > 0; idx -= idx & -idx) {
            res += tree[idx];
        }
        return res;
    }

    long long rangeSum(int l, int r) {
        return sumPrefix(r) - sumPrefix(l - 1);
    }
};
\end{lstlisting}

\subsection{适用场景}

单点修改、区间和查询。

\section{T402 线段树}

\subsection{适用场景}

区间加、区间和查询。模板较长，只有练过才建议考场使用。

\begin{lstlisting}
struct SegTree {
    int n;
    vector<long long> tree, lazy;

    SegTree(int n) : n(n), tree(4 * n + 4), lazy(4 * n + 4) {}

    void apply(int node, int l, int r, long long val) {
        tree[node] += val * (r - l + 1);
        lazy[node] += val;
    }

    void push(int node, int l, int r) {
        if (lazy[node] == 0 || l == r) return;
        int mid = (l + r) / 2;
        apply(node * 2, l, mid, lazy[node]);
        apply(node * 2 + 1, mid + 1, r, lazy[node]);
        lazy[node] = 0;
    }

    void rangeAdd(int node, int l, int r, int ql, int qr, long long val) {
        if (ql <= l && r <= qr) {
            apply(node, l, r, val);
            return;
        }
        push(node, l, r);
        int mid = (l + r) / 2;
        if (ql <= mid) rangeAdd(node * 2, l, mid, ql, qr, val);
        if (qr > mid) rangeAdd(node * 2 + 1, mid + 1, r, ql, qr, val);
        tree[node] = tree[node * 2] + tree[node * 2 + 1];
    }

    long long query(int node, int l, int r, int ql, int qr) {
        if (ql <= l && r <= qr) return tree[node];
        push(node, l, r);
        int mid = (l + r) / 2;
        long long res = 0;
        if (ql <= mid) res += query(node * 2, l, mid, ql, qr);
        if (qr > mid) res += query(node * 2 + 1, mid + 1, r, ql, qr);
        return res;
    }
};
\end{lstlisting}

调用：

\begin{lstlisting}
SegTree seg(n);
seg.rangeAdd(1, 1, n, l, r, v);
long long ans = seg.query(1, 1, n, l, r);
\end{lstlisting}

\subsection{易错点}

\begin{itemize}
  \item 默认区间是 1 到 $n$。
  \item 懒标记下传不要漏。
  \item 求和用 \Code{long long}。
\end{itemize}

\section{T403 单调栈}

\subsection{适用场景}

找每个元素右边第一个更大元素、左边第一个更小元素等。

\begin{lstlisting}
vector<int> ans(n, -1);
stack<int> st; // indices

for (int i = 0; i < n; i++) {
    while (!st.empty() && a[i] > a[st.top()]) {
        ans[st.top()] = i;
        st.pop();
    }
    st.push(i);
}
\end{lstlisting}

\subsection{易错点}

\begin{itemize}
  \item 比较符号决定严格更大还是大于等于。
  \item 栈里通常存下标，不是值。
  \item 全相等数据要测试。
\end{itemize}

\section{T404 单调队列}

\subsection{适用场景}

固定长度滑动窗口最大值或最小值。

\begin{lstlisting}
deque<int> dq; // indices, values decreasing
for (int i = 0; i < n; i++) {
    while (!dq.empty() && dq.front() <= i - k) dq.pop_front();
    while (!dq.empty() && a[dq.back()] <= a[i]) dq.pop_back();
    dq.push_back(i);

    if (i >= k - 1) {
        cout << a[dq.front()] << '\n';
    }
}
\end{lstlisting}

\subsection{易错点}

\begin{itemize}
  \item 先删除过期下标。
  \item 最大值队列保持递减，最小值队列保持递增。
  \item 输出窗口从 \Code{i >= k - 1} 开始。
\end{itemize}

\section{T501 split 字符串分割}

\subsection{按指定字符分割}

\begin{lstlisting}
vector<string> split(const string& s, char sep) {
    vector<string> res;
    string cur;
    for (char c : s) {
        if (c == sep) {
            res.push_back(cur);
            cur.clear();
        } else {
            cur.push_back(c);
        }
    }
    res.push_back(cur);
    return res;
}
\end{lstlisting}

\subsection{按空格分割}

\begin{lstlisting}
stringstream ss(line);
string word;
while (ss >> word) {
    // use word
}
\end{lstlisting}

\section{T502 KMP}

\begin{lstlisting}
vector<int> prefixFunction(const string& p) {
    int n = (int)p.size();
    vector<int> pi(n, 0);
    for (int i = 1; i < n; i++) {
        int j = pi[i - 1];
        while (j > 0 && p[i] != p[j]) j = pi[j - 1];
        if (p[i] == p[j]) j++;
        pi[i] = j;
    }
    return pi;
}

int countMatch(const string& text, const string& p) {
    vector<int> pi = prefixFunction(p);
    int j = 0, cnt = 0;
    for (char c : text) {
        while (j > 0 && c != p[j]) j = pi[j - 1];
        if (c == p[j]) j++;
        if (j == (int)p.size()) {
            cnt++;
            j = pi[j - 1];
        }
    }
    return cnt;
}
\end{lstlisting}

\section{T503 Trie}

\begin{lstlisting}
struct Trie {
    struct Node {
        int child[26];
        bool end = false;
        Node() {
            memset(child, -1, sizeof child);
        }
    };

    vector<Node> tr;

    Trie() {
        tr.push_back(Node());
    }

    void insert(const string& s) {
        int u = 0;
        for (char c : s) {
            int x = c - 'a';
            if (tr[u].child[x] == -1) {
                tr[u].child[x] = (int)tr.size();
                tr.push_back(Node());
            }
            u = tr[u].child[x];
        }
        tr[u].end = true;
    }

    bool find(const string& s) {
        int u = 0;
        for (char c : s) {
            int x = c - 'a';
            if (tr[u].child[x] == -1) return false;
            u = tr[u].child[x];
        }
        return tr[u].end;
    }
};
\end{lstlisting}

\Warn{这个 Trie 模板只适合小写英文字母。}

\section{T510 整除取整和标准取模}

\subsection{适用场景}

\begin{itemize}
  \item 分组、分页、分块，要求正整数向上取整。
  \item 坐标、时间差、偏移量可能为负，要求数学意义的向上或向下取整。
  \item 同余、周期、哈希下标中需要把负数余数转成 $0$ 到 $m-1$。
\end{itemize}

\subsection{正整数向上取整}

若 $a \ge 0, b > 0$：

\begin{lstlisting}
long long ceil_div_positive(long long a, long long b) {
    return (a + b - 1) / b;
}
\end{lstlisting}

示例：

\begin{lstlisting}
long long pages = ceil_div_positive(n, pageSize);
long long blocks = ceil_div_positive(n, blockSize);
\end{lstlisting}

\subsection{有符号通用版本}

\begin{lstlisting}
long long floor_div(long long a, long long b) {
    long long q = a / b;
    long long r = a % b;
    if (r != 0 && ((r > 0) != (b > 0))) q--;
    return q;
}

long long ceil_div(long long a, long long b) {
    long long q = a / b;
    long long r = a % b;
    if (r != 0 && ((r > 0) == (b > 0))) q++;
    return q;
}
\end{lstlisting}

\subsection{标准非负余数}

\begin{lstlisting}
long long norm_mod(long long x, long long m) {
    long long r = x % m;
    if (r < 0) r += m;
    return r;
}
\end{lstlisting}

\subsection{易错点}

\begin{itemize}
  \item \Code{(a+b-1)/b} 只适合 $a \ge 0, b > 0$，不是通用向上取整。
  \item C++ 整数除法向 0 截断，\Code{-3 / 2} 是 \Code{-1}。
  \item C++ 负数取模可能得到负数，\Code{-1 \% 5} 是 \Code{-1}。
  \item 除数 \Code{b} 或模数 \Code{m} 不能为 0。
  \item 若 \Code{a + b - 1} 可能溢出，需要改写或先确认范围。
\end{itemize}

\section{T511 快速幂}

\subsection{适用场景}

计算 $a^b \bmod mod$。

\begin{lstlisting}
long long modPow(long long a, long long b, long long mod) {
    long long res = 1 % mod;
    a %= mod;
    while (b > 0) {
        if (b & 1) res = res * a % mod;
        a = a * a % mod;
        b >>= 1;
    }
    return res;
}
\end{lstlisting}

\section{T512 素数和约数}

\subsection{gcd 和 lcm}

\begin{lstlisting}
long long g = gcd(a, b);
long long l = a / g * b;
\end{lstlisting}

\subsection{判断素数}

\begin{lstlisting}
bool isPrime(long long x) {
    if (x < 2) return false;
    for (long long d = 2; d * d <= x; d++) {
        if (x % d == 0) return false;
    }
    return true;
}
\end{lstlisting}

\subsection{枚举约数}

\begin{lstlisting}
vector<long long> divisors;
for (long long d = 1; d * d <= x; d++) {
    if (x % d == 0) {
        divisors.push_back(d);
        if (d != x / d) divisors.push_back(x / d);
    }
}
sort(divisors.begin(), divisors.end());
\end{lstlisting}

\section{T513 进制转换}

\subsection{任意进制转十进制}

\begin{lstlisting}
long long toDecimal(const string& s, int base) {
    long long ans = 0;
    for (char c : s) {
        int v;
        if ('0' <= c && c <= '9') v = c - '0';
        else if ('A' <= c && c <= 'F') v = c - 'A' + 10;
        else v = c - 'a' + 10;
        ans = ans * base + v;
    }
    return ans;
}
\end{lstlisting}

\subsection{十进制转任意进制}

\begin{lstlisting}
string fromDecimal(long long x, int base) {
    if (x == 0) return "0";
    string digits = "0123456789ABCDEF";
    string res;
    while (x > 0) {
        res.push_back(digits[x % base]);
        x /= base;
    }
    reverse(res.begin(), res.end());
    return res;
}
\end{lstlisting}

\Warn{数字很长时不能转成 \Code{long long}，要用字符串模拟。}

\section{T514 日期处理}

\subsection{闰年}

\begin{lstlisting}
bool leap(int y) {
    return (y % 400 == 0) || (y % 4 == 0 && y % 100 != 0);
}
\end{lstlisting}

\subsection{每月天数}

\begin{lstlisting}
int mdays(int y, int m) {
    int d[] = {0,31,28,31,30,31,30,31,31,30,31,30,31};
    if (m == 2 && leap(y)) return 29;
    return d[m];
}
\end{lstlisting}

\subsection{下一天}

\begin{lstlisting}
void nextDay(int& y, int& m, int& d) {
    d++;
    if (d > mdays(y, m)) {
        d = 1;
        m++;
        if (m > 12) {
            m = 1;
            y++;
        }
    }
}
\end{lstlisting}

\subsection{时间转秒}

\begin{lstlisting}
int toSeconds(int h, int m, int s) {
    return h * 3600 + m * 60 + s;
}
\end{lstlisting}

\section{T601 大模拟骨架}

\subsection{适用场景}

C 题长题面、多对象、多命令、多状态。

\begin{lstlisting}
struct Item {
    int id;
    string name;
    int state;
};

int n, q;
vector<Item> items;

void readInput() {
    cin >> n >> q;
    items.resize(n);
    for (int i = 0; i < n; i++) {
        cin >> items[i].id >> items[i].name;
        items[i].state = 0;
    }
}

void solve() {
    while (q--) {
        string op;
        cin >> op;
        // dispatch command
    }
}

int main() {
    ios::sync_with_stdio(false);
    cin.tie(nullptr);

    readInput();
    solve();
    return 0;
}
\end{lstlisting}

\section{T602 命令分发}

\begin{lstlisting}
void handleAdd() {
    int id, x;
    cin >> id >> x;
    // add logic
}

void handleDelete() {
    int id;
    cin >> id;
    // delete logic
}

void handleQuery() {
    int id;
    cin >> id;
    // query logic
}

void handleCommand() {
    string op;
    cin >> op;
    if (op == "add") handleAdd();
    else if (op == "delete") handleDelete();
    else if (op == "query") handleQuery();
}
\end{lstlisting}

\subsection{易错点}

\begin{itemize}
  \item 每种命令的参数个数不同。
  \item 命令参数可能包含空格时要用 \Code{getline}。
  \item 查询命令通常直接输出，修改命令通常不输出。
\end{itemize}

\section{T603 状态机}

\begin{lstlisting}
const int WAITING = 0;
const int RUNNING = 1;
const int FINISHED = 2;

struct Job {
    int id;
    int status;
    int score;
};

void start(Job& job) {
    if (job.status != WAITING) return;
    job.status = RUNNING;
}

void finish(Job& job) {
    if (job.status != RUNNING) return;
    job.status = FINISHED;
    job.score += 10;
}
\end{lstlisting}

\subsection{检查点}

\begin{itemize}
  \item 初始状态是什么。
  \item 每个操作允许在哪些状态发生。
  \item 非法状态转换是忽略还是输出错误。
  \item 状态变化时是否要同步修改其他字段。
\end{itemize}

\section{T604 懒删除优先队列}

\subsection{适用场景}

动态更新或删除元素，同时需要不断取当前最优元素。

\begin{lstlisting}
struct Task {
    int priority;
    int id;
};

struct Cmp {
    bool operator()(const Task& a, const Task& b) const {
        if (a.priority != b.priority) return a.priority < b.priority;
        return a.id > b.id;
    }
};

priority_queue<Task, vector<Task>, Cmp> pq;
map<int, int> currentPriority;

void addOrUpdate(int id, int priority) {
    currentPriority[id] = priority;
    pq.push({priority, id});
}

void removeTask(int id) {
    currentPriority.erase(id);
}

Task getTop() {
    while (!pq.empty()) {
        Task t = pq.top();
        if (currentPriority.count(t.id) &&
            currentPriority[t.id] == t.priority) {
            return t;
        }
        pq.pop();
    }
    return {-1, -1};
}
\end{lstlisting}

\subsection{易错点}

\begin{itemize}
  \item 每次取堆顶前都要清理过期元素。
  \item 更新元素时直接压入新值，不删除旧值。
  \item 判断有效性必须和当前状态表一致。
\end{itemize}

\section{T605 LRU 缓存}

\subsection{适用场景}

最近最少使用淘汰，访问后变成最新。

\begin{lstlisting}
int cap;
list<int> order; // front is most recent
unordered_map<int, list<int>::iterator> pos;

void access(int x) {
    if (cap == 0) return;

    if (pos.count(x)) {
        order.erase(pos[x]);
    } else if ((int)order.size() == cap) {
        int old = order.back();
        order.pop_back();
        pos.erase(old);
    }

    order.push_front(x);
    pos[x] = order.begin();
}
\end{lstlisting}

\subsection{易错点}

\begin{itemize}
  \item 容量为 0 要特判。
  \item 删除链表节点后，map 里的迭代器也要更新。
  \item \Code{front} 是最新还是最旧要统一。
\end{itemize}

\section{T606 路径规范化}

\subsection{适用场景}

文件路径、目录路径，处理 \Code{.}、\Code{..}、多余斜杠。

\begin{lstlisting}
vector<string> normalizePath(const string& path) {
    vector<string> st;
    vector<string> parts = split(path, '/');
    for (string part : parts) {
        if (part.empty() || part == ".") {
            continue;
        } else if (part == "..") {
            if (!st.empty()) st.pop_back();
        } else {
            st.push_back(part);
        }
    }
    return st;
}

string joinPath(const vector<string>& parts) {
    if (parts.empty()) return "/";
    string res;
    for (const string& p : parts) {
        res += "/";
        res += p;
    }
    return res;
}
\end{lstlisting}

\Warn{如果题目允许相对路径，是否能越过根目录要按题意处理。}

\section{T607 URL 参数解析}

\subsection{适用场景}

形如 \Code{/path/to?a=1\&b=2} 的字符串。

\begin{lstlisting}
struct Url {
    string path;
    map<string, string> query;
};

Url parseUrl(const string& s) {
    Url res;
    int qpos = (int)s.find('?');
    if (qpos == (int)string::npos) {
        res.path = s;
        return res;
    }

    res.path = s.substr(0, qpos);
    string qs = s.substr(qpos + 1);
    vector<string> pairs = split(qs, '&');
    for (string p : pairs) {
        int eq = (int)p.find('=');
        if (eq == (int)string::npos) {
            res.query[p] = "";
        } else {
            string key = p.substr(0, eq);
            string val = p.substr(eq + 1);
            res.query[key] = val;
        }
    }
    return res;
}
\end{lstlisting}

\subsection{易错点}

\begin{itemize}
  \item URL 可能没有 \Code{?}。
  \item 参数值可能为空。
  \item 重复参数时保留第一个还是最后一个要看题目。
\end{itemize}

\section{T608 括号匹配}

\subsection{适用场景}

判断括号是否合法，或处理简单嵌套结构。

\begin{lstlisting}
bool checkBrackets(const string& s) {
    stack<char> st;
    for (char c : s) {
        if (c == '(' || c == '[' || c == '{') {
            st.push(c);
        } else if (c == ')' || c == ']' || c == '}') {
            if (st.empty()) return false;
            char t = st.top();
            st.pop();
            if (c == ')' && t != '(') return false;
            if (c == ']' && t != '[') return false;
            if (c == '}' && t != '{') return false;
        }
    }
    return st.empty();
}
\end{lstlisting}

\section{T609 简单表达式求值}

\subsection{只含加减和非负整数}

\begin{lstlisting}
long long evalPlusMinus(const string& s) {
    long long ans = 0;
    long long num = 0;
    int sign = 1;

    for (int i = 0; i <= (int)s.size(); i++) {
        if (i < (int)s.size() && isdigit(s[i])) {
            num = num * 10 + (s[i] - '0');
        } else {
            ans += sign * num;
            num = 0;
            if (i < (int)s.size()) {
                if (s[i] == '+') sign = 1;
                else if (s[i] == '-') sign = -1;
            }
        }
    }
    return ans;
}
\end{lstlisting}

\subsection{易错点}

\begin{itemize}
  \item 上面模板不处理乘除和括号。
  \item 如果有空格，先跳过空格。
  \item 如果有负数开头，检查是否符合模板逻辑。
\end{itemize}
